\chapter{Le métier de teneur de marché}\label{ch:hf:the-business-of-market-making}

En 2025, Virtu Financial, un \omterm{def:hf:the-business-of-market-making:emm}{teneur de marché électronique} coté, a publié un revenu net ajusté de négociation de 8.6~millions de dollars par jour de négociation. Le même dépôt montre
ce que ce revenu a coûté~: 770~millions de dollars de frais de bourse, de compensation et de flux d'ordres, 528~millions de dollars de rémunérations pour environ mille employés, 249~millions de dollars de
communication et de données. Flow Traders, spécialiste des fonds indiciels cotés, a gagné 2.5 points de base sur chaque euro des produits qu'il a négociés. Le
premier teneur de marché de ce chapitre, exploité pendant dix heures simulées, gagne un demi-centime par action sur la fourchette et restitue presque tout en dix secondes. Le
métier repose sur une faible marge appliquée à un très gros volume, sur une base de coûts qui ne diminue pas quand le volume baisse.

\section{Ce que vend un teneur de marché}

Un teneur de marché (livre~1, chapitre~1) se tient prêt à acheter à son cours acheteur et à vendre à son cours vendeur. Ce qu'il vend, c'est l'immédiateté~: l'investisseur qui
traverse la fourchette négocie maintenant au lieu d'attendre quelqu'un qui veut la contrepartie. Le prix de l'immédiateté est la moitié de la fourchette par action, gagnée
sur chaque exécution, et son coût tient à ce que certains de ceux qui la prennent savent où va le prix.

\begin{definition}[Teneur de marché électronique]\label{def:hf:the-business-of-market-making:emm}
Un \emph{teneur de marché électronique}\index{teneur de marché électronique} (electronic market maker) est une entreprise qui fournit de la liquidité par algorithme~: ses cotations, leurs tailles et leurs
révisions sont fixées par des programmes qui réagissent aux données de marché en microsecondes à millisecondes, sur de nombreux instruments et plateformes à la fois, des humains
supervisant les programmes plutôt que les cotations.
\end{definition}

La plupart des \omterm{def:hf:the-business-of-market-making:emm}{teneurs de marché électroniques} sont aussi des traders à haute fréquence au sens juridique que le livre~10, chapitre~24 donne au terme (durées de
détention courtes, taux de messages élevés, serveurs colocalisés)~; la réciproque n'est pas vraie, car certaines stratégies à haute fréquence se bornent à prendre de la liquidité. Certains \omterm{def:hf:the-business-of-market-making:emm}{teneurs de marché électroniques}
portent des obligations, en tant que teneurs de marché désignés ou teneurs de marché principaux (livre~1, chapitres~19 et~29)~; la plupart cotent parce que cela rapporte. Quatre lignes d'activité
reviennent dans ce livre~: les produits négociés en bourse et leurs couvertures (chapitres~12 à~14), les contrats à terme et les options sur bourses (chapitres~15, 19 et~20), les
marchés de dealers atteints par demande de cotation (RFQ) ou par cotation en continu, change et obligations (chapitres~21 et~22), et le flux d'ordres de détail acheté auprès de courtiers (chapitre~23).
Les plateformes crypto (chapitres~24 et~25) et les marchés d'événements (chapitre~26) font la même opération sur une infrastructure plus récente.

\section{Revenu~: capture, volume et volatilité}

\begin{definition}[Taux de capture du revenu]\label{def:hf:the-business-of-market-making:capture}
Le \emph{taux de capture du revenu}\index{taux de capture du revenu} (revenue capture) d'un teneur de marché est son revenu de négociation, net des frais qui varient avec les transactions, par unité négociée~: par
action, par contrat, ou par unité de valeur négociée, généralement en points de base.
\end{definition}

Le revenu est égal au taux de capture multiplié par le volume. Les deux évoluent avec la volatilité~: des journées plus volatiles apportent plus de volume, des fourchettes plus larges et des mouvements défavorables plus importants, et les
entreprises cotées décrivent leur revenu comme dépendant à la fois du volume et de la volatilité. La capture elle-même est ce qui reste de la demi-fourchette une fois que le prix
a bougé contre les exécutions. Le livre~7, chapitre~23 la décompose précisément.

\begin{proposition}[Où va la demi-fourchette]\label{prop:hf:the-business-of-market-making:split}
Soient $i$ les exécutions de côté $s_i=\pm1$, de quantité $q_i$, de prix $p_i$ et d'instant $t_i$~; soit $m(t)$ un prix de référence, $H$ un horizon et $T$ la fin de
la période~; soit $f_i$ les frais payés sur l'exécution $i$ (négatifs pour une rétrocession). Le P\&L des exécutions, la position finale valorisée en $m(T)$, vaut
\begin{align*}
\pnl={}&\underbrace{\sum_i s_iq_i\,(m(t_i)-p_i)}_{\text{capture du spread}}
+\underbrace{\sum_i s_iq_i\,\bigl(m(t_i+H)-m(t_i)\bigr)}_{\text{sélection adverse}}\\
&+\underbrace{\sum_i s_iq_i\,\bigl(m(T)-m(t_i+H)\bigr)}_{\text{inventaire}}-\sum_i f_i .
\end{align*}
\end{proposition}

\begin{proof}
Pour chaque exécution, les trois crochets se télescopent en $m(T)-p_i$~; $\sum_i s_iq_i\,(m(T)-p_i)$ est la trésorerie issue des exécutions augmentée de la position finale valorisée en
$m(T)$.
\end{proof}

L'identité vaut pour toute référence~: le mid, un prix juste (chapitre~2), ou, dans une simulation, le prix efficient $P^\ast_t$ que personne n'observe. Avec
le mid, la capture du spread d'une exécution passive vaut exactement la moitié de la fourchette~; le terme de sélection adverse mesure la part que les $H$ secondes suivantes en
reprennent. L'étude de Menkveld sur un grand teneur de marché à haute fréquence, sur Chi-X et Euronext, retrouve la même décomposition dans des données réelles~: un profit brut de
0.88~\euro par transaction, formé de 1.55~\euro gagnés sur la fourchette, nets de frais, et d'une perte de 0.68~\euro sur les positions, qui ont rapporté 0.45~\euro lorsqu'elles étaient
détenues moins de cinq secondes et ont perdu 1.13~\euro lorsqu'elles étaient détenues plus longtemps.

\section{Coûts~: frais, financement, personnel et machines}

Les coûts d'un teneur de marché se répartissent en deux groupes. Ceux qui varient avec les transactions sont les frais de bourse et de compensation (diminués des rétrocessions), le paiement pour flux d'ordres
lorsqu'il achète des ordres de détail, et le financement de son inventaire ainsi que de la marge qui le couvre. Ceux qui ne varient pas sont le personnel, les données de marché,
la connectivité et la colocation, les serveurs et les adhésions aux bourses. Le premier groupe est déduit de la capture~; le second détermine la capture minimale dont
l'entreprise a besoin pour survivre.

\begin{omfigure}
\begin{tikzpicture}
\begin{axis}[width=12cm,height=5.6cm,ybar stacked,bar width=20pt,ymin=0,ymax=4000,ylabel={\small millions de dollars},
  xtick={0,1,2,3,4,5,6,7},xticklabels={revenu,frais,intérêts,personnel,données,{opé\-rations},{amor\-tissements},{ce qui reste}},
  xticklabel style={font=\scriptsize,text width=1.5cm,align=center},tick label style={font=\footnotesize},
  table/col sep=comma,enlarge x limits=0.08,grid=major,grid style={gray!20},axis x line*=bottom]
\addplot[draw=none,fill=none] table[x=k,y=bottom]{figdata/hft/01-the-business-of-market-making/waterfall.csv};
\addplot[fill=omThm!50,draw=omThm] table[x=k,y=cost]{figdata/hft/01-the-business-of-market-making/waterfall.csv};
\addplot[fill=omDef!60,draw=omDef] table[x=k,y=total]{figdata/hft/01-the-business-of-market-making/waterfall.csv};
\end{axis}
\end{tikzpicture}

\omcaption{Revenu de Virtu Financial en 2025 et coûts qu'il a supportés jusqu'au bas du graphique, en millions de dollars~: les frais de bourse, de compensation et de flux d'ordres et
les charges d'intérêts et de dividendes varient avec les transactions~; le personnel, la communication et les données, les opérations et les amortissements n'en dépendent pas. Ce qui reste est avant intérêts de financement, éléments exceptionnels
et impôts. Données~: formulaire 10-K 2025.}\label{fig:hf:the-business-of-market-making:waterfall}
\end{omfigure}

\begin{dated}{2026-09}{Deux teneurs de marché cotés en 2025}\label{dat:hf:the-business-of-market-making:public}
\textbf{Virtu Financial} (formulaire 10-K 2025)~: revenu total 3\,632~millions de dollars, dont revenu de négociation 2\,437~millions de dollars~; frais de courtage, de bourse,
de compensation et paiements pour flux d'ordres 770~millions de dollars~; charges d'intérêts et de dividendes 647~millions de dollars~; rémunération des employés 528~millions de dollars~;
communication et traitement des données 249~millions de dollars. Revenu net ajusté de négociation (la mesure de l'entreprise~: revenu de négociation, commissions et produits d'intérêts et de
dividendes, diminués des deux lignes de coûts directs) 2\,145~millions de dollars, soit 8.6~millions de dollars par jour sur 248.5 jours de négociation~; le segment teneur de marché
1\,666~millions de dollars, soit 6.7~millions de dollars par jour, 77.7~\% du total. Environ 1\,027 employés en février 2026.

\textbf{Flow Traders} (résultats annuels 2025)~: revenu net de négociation 485.8~millions d'euros sur 1\,940~milliards d'euros de valeur négociée en produits négociés en bourse~;
charges opérationnelles fixes 204.1~millions d'euros (personnel fixe 97.3~millions d'euros, technologie 70.6~millions d'euros, autres 36.3~millions d'euros), charges variables
de personnel 77.4~millions d'euros~; EBITDA 198.9~millions d'euros, soit une marge de 41~\%~; 635 employés à temps plein à la fin de l'année.
\end{dated}

Trois ratios se dégagent de ces données. Les lignes de coûts fixes de Virtu (personnel, communication et données, opérations, amortissements) totalisent 940~millions de dollars, soit 3.8
millions de dollars par jour de négociation, contre 8.6~millions de dollars de revenu net de négociation~; la seule communication et les données coûtent 1.0~million de dollars par jour de négociation. Son revenu net ajusté
de négociation s'élevait à 2.1~millions de dollars par employé, contre 0.51~millions de dollars de rémunération. Le revenu net de négociation de Flow Traders représentait 2.5 points de base de
la valeur de ses produits négociés en bourse~; le ratio surestime ce que l'entreprise a gagné sur ces produits, car le revenu inclut aussi d'autres activités de négociation. Ses
coûts fixes atteignaient 0.32~millions d'euros par employé, dont un tiers pour la technologie.

\section{Ce que la littérature publique montre sur les gagnants}

Les deux dépôts décrivent des entreprises, non des transactions. La littérature académique ajoute ce que rapporte une transaction et qui en profite. La décomposition de Menkveld établit que la
fourchette rapporte et que la position coûte, de sorte que le métier consiste à gagner la première et à raccourcir la seconde. Baron, Brogaard, Hagströmer et Kirilenko ont classé
les traders à haute fréquence selon leur latence relative et ont montré que les écarts de vitesse expliquent de larges écarts de performance~; les entreprises qui ont progressé au
classement grâce à une amélioration de leur colocation ont ensuite gagné davantage, aussi bien grâce à l'information de courte durée qu'à une gestion plus rapide des risques, sur le market making comme
dans l'arbitrage entre marchés. La vitesse n'est donc pas un détail de mise en œuvre mais une partie du produit~: elle détermine à quelle fréquence la cotation d'un teneur de marché
est celle qu'un trader mieux informé saisit (chapitre~9), et à quelle vitesse il peut se couvrir.

\begin{remark}[Ce qu'un dépôt ne permet pas de dire]\label{rem:hf:the-business-of-market-making:limits}
Un dépôt agrège des milliers d'instruments et de nombreuses stratégies. Il ne donne aucun taux de capture par action pour l'entreprise dans son ensemble (Virtu ne communique pas son
volume), aucune répartition entre négociation passive et négociation agressive, et aucune vue des journées perdantes au sein d'une année. Le reste de ce livre travaille à partir de
simulations et de la recherche publiée pour ces raisons précises, et le dit dans chaque légende.
\end{remark}

\section{La structure de l'activité~: coûts fixes et échelle}

Avec un taux de capture $c$ par unité échangée, des coûts variables $v$ par unité (frais, financement) et des coûts fixes $F$ par jour, un volume quotidien $V$ rapporte
\[
\pi(V)=(c-v)\,V-F .
\]

\begin{proposition}[Volume du seuil de rentabilité et levier opérationnel]\label{prop:hf:the-business-of-market-making:leverage}
Si $c>v$, le profit est nul en $V^\ast=F/(c-v)$, et au-delà l'élasticité du profit au volume vaut
\[
\frac{d\ln\pi}{d\ln V}=\frac{(c-v)\,V}{(c-v)\,V-F}=\frac{V}{V-V^\ast}>1 .
\]
\end{proposition}

\begin{proof}
Dérivation directe de $\pi(V)$~; la seconde forme divise le numérateur et le dénominateur par $c-v$.
\end{proof}

\begin{example}[Une petite entreprise]\label{ex:hf:the-business-of-market-making:small}
Un nouveau teneur de marché attend un taux de capture brut de 0.5 point de base de valeur échangée, paie 0.2 point de base en frais et financement, et coûte 40\,000~\$ par
jour de négociation à faire tourner (environ 10~millions de dollars par an). Il atteint le seuil de rentabilité à 1.33~milliard de dollars de valeur échangée par jour. À 2~milliards de dollars, il gagne 20\,000~\$ par jour,
avec un levier opérationnel de 3~: une baisse de 10~\% du volume retire 30~\% de son profit. Si la concurrence fait passer le taux de capture de 0.5 à 0.4 point de base, le
volume du seuil de rentabilité monte à 2~milliards de dollars, et la même entreprise ne gagne rien (\cref{fig:hf:the-business-of-market-making:breakeven}).
\end{example}

\begin{omfigure}
\begin{tikzpicture}
\begin{axis}[width=11cm,height=5.2cm,xlabel={\small valeur quotidienne échangée (milliards de dollars)},ylabel={\small profit (milliers de dollars par jour)},
  tick label style={font=\footnotesize},xmin=0,xmax=3,ymin=-45,ymax=55,grid=major,grid style={gray!25},table/col sep=comma]
\addplot[omThm,thick] table[x=volume_bn,y=profit_k]{figdata/hft/01-the-business-of-market-making/breakeven.csv};
\draw[gray,dashed] (axis cs:0,0) -- (axis cs:3,0);
\draw[omDef,thick] (axis cs:1.333,-45) -- (axis cs:1.333,0);
\node[font=\scriptsize,anchor=west,omDef,fill=white,inner sep=1pt] at (axis cs:1.4,-30) {seuil de rentabilité~: 1.33~Md~\$};
\end{axis}
\end{tikzpicture}

\omcaption{Profit quotidien de la petite entreprise de l'\cref{ex:hf:the-business-of-market-making:small} (capture 0.5 point de base, coûts variables 0.2,
coûts fixes 40\,000~\$ par jour) en fonction de la valeur quotidienne échangée. Données~: \texttt{firm.mmecon}.}\label{fig:hf:the-business-of-market-making:breakeven}
\end{omfigure}

La base de coûts fixes explique trois caractéristiques du secteur que les chapitres suivants reprennent. Les teneurs de marché répartissent une même pile technologique sur autant
d'instruments et de plateformes qu'elle peut en servir, car chaque nouvel instrument ajoute du volume pour un coût fixe faible. Ils paient pour la vitesse, dont le coût est fixe
et dont le bénéfice croît avec le volume. Et les entreprises qui survivent à une longue période de faible volatilité et de faible volume sont celles qui ont le volume de seuil de rentabilité
le plus bas, et non le taux de capture le plus élevé.

\section{Tutoriel~: un premier teneur de marché}\label{tut:hf:the-business-of-market-making:tut}

\begin{tutorial}
\textbf{Objectif.} Faire tourner le teneur de marché le plus simple sur un marché simulé et décomposer son P\&L. \textbf{État final~:} le tableau ci-dessous et la
\cref{fig:hf:the-business-of-market-making:markout}.
\begin{tutsteps}
\item \textbf{Le marché.} \texttt{firm.tape} du livre~7 simule une action à 100~\$ avec un tick d'un cent pendant une heure~: des fournisseurs de liquidité, des noise
traders dont les ordres se regroupent, des traders informés qui échangent vers un prix efficient $P^\ast_t$ qui saute à des instants aléatoires, et une fenêtre d'annonce.
\item \textbf{Le harnais.} \texttt{firm.mmharness} fait tourner un \emph{Quoter} à l'intérieur de ce marché~: les ordres du Quoter sont de vrais ordres dans le carnet, il
voit le marché avec un retard d'une demi-milliseconde et ses ordres arrivent une demi-milliseconde après leur envoi. Il doit coter à partir du carnet \emph{sans}
ses propres ordres, sinon il les poursuit~: une version qui s'alignait sur le meilleur cours acheteur brut, ses propres ordres inclus, a envoyé 336\,012 messages dans les
premières 1\,400 secondes simulées d'une même exécution, presque tous dans sa dernière minute, contre 323 pour la version corrigée. \texttt{Ctx.quote} conserve un ordre par côté à un prix cible
(\cref{lst:hf:the-business-of-market-making:quote}).
\omcode{code/firm/mmharness/firm_mmharness.py}{135}{147}{Conserver un ordre par côté au prix cible~: annuler les autres, en envoyer un si aucun ne reste.}\label{lst:hf:the-business-of-market-making:quote}
\item \textbf{Le Quoter.} \texttt{SymmetricQuoter} s'aligne sur le meilleur cours acheteur et le meilleur cours vendeur des autres avec 100 actions et cesse d'ajouter à une position de
500 (\cref{lst:hf:the-business-of-market-making:quoter}). Les frais sont illustratifs~: une rémunération de 0.1 cent par action sur les exécutions passives, un prélèvement de 0.1 cent
sur les exécutions agressives.
\omcode{code/firm/mmharness/firm_mmharness.py}{166}{184}{Le premier teneur de marché du chapitre.}\label{lst:hf:the-business-of-market-making:quoter}
\item \textbf{Exécuter et décomposer.} \texttt{hf\_business.decomposition(H=10)} fait tourner dix heures indépendantes (graines 1--10) et répartit le P\&L comme dans la
\cref{prop:hf:the-business-of-market-making:split}, par rapport au mid et par rapport au prix efficient.
\end{tutsteps}
\end{tutorial}

Le Quoter échange en moyenne 38\,270 actions par heure, soit 11~\% du volume du marché, et envoie 1\,256 messages. Son P\&L par heure et par action (moyennes
sur dix heures, $H=10$ secondes)~:

\begin{center}
\small
\begin{tabular}{@{}lrrrrr@{}}
\toprule
référence & capture du spread & sélection adverse & inventaire & frais & total\\
\midrule
mid, \$ par heure & 186.50 & $-177.25$ & $-66.50$ & 38.27 & $-18.98$\\
mid, cents par action & 0.487 & $-0.463$ & $-0.174$ & 0.100 & $-0.050$\\
prix efficient, cents par action & $-0.173$ & 0.038 & $-0.013$ & 0.100 & $-0.049$\\
\bottomrule
\end{tabular}
\end{center}

Par rapport au mid, le Quoter gagne sa demi-fourchette, un demi-cent par action, et la perd presque entièrement au profit de la sélection adverse en moins de dix secondes~;
la rémunération est ce qui le maintient près de zéro. Par rapport au prix efficient, la capture du spread est déjà négative~: en moyenne, les exécutions du Quoter ont lieu
lorsque le prix efficient a déjà traversé sa cotation, de sorte que le mid auquel il échange est périmé. La courbe de mark-out montre la demi-fourchette qui fond
(\cref{fig:hf:the-business-of-market-making:markout}). Le total horaire a un écart type de 89~\$, si bien que la moyenne sur dix heures, de $-\$19$, porte une
erreur type de 28~\$~: le Quoter naïf ne gagne pas d'argent, et dix heures ne permettent pas de dire s'il en perd.

\begin{omfigure}
\begin{tikzpicture}
\begin{axis}[width=11cm,height=5.2cm,xlabel={\small secondes après l'exécution},ylabel={\small mark-out (cents par action)},xmode=log,log basis x=10,
  xmin=0.4,xmax=70,ymin=-0.3,ymax=0.55,tick label style={font=\footnotesize},grid=major,grid style={gray!25},table/col sep=comma,
  xtick={0.5,1,2,5,10,20,60},xticklabels={0.5,1,2,5,10,20,60},legend style={font=\scriptsize,at={(0.5,-0.3)},anchor=north},legend columns=2]
\addplot[omDef,thick,mark=*,mark size=1.2pt] table[x=h,y=mid]{figdata/hft/01-the-business-of-market-making/markout.csv};
\addplot[omThm,thick,dashed,mark=square*,mark size=1.2pt] table[x=h,y=truth]{figdata/hft/01-the-business-of-market-making/markout.csv};
\draw[gray] (axis cs:0.4,0) -- (axis cs:70,0);
\legend{par rapport au mid,par rapport au prix efficient}
\end{axis}
\end{tikzpicture}

\omcaption{Mark-out des exécutions du Quoter symétrique, $s_i\,(m(t_i+h)-p_i)$ en cents par action, en fonction de l'horizon $h$ (échelle logarithmique), avec $m$ le mid ou
le prix efficient~; pondéré par les quantités sur dix heures simulées. En $h=0$, non tracé sur l'échelle logarithmique, le mark-out par rapport au mid est la demi-fourchette, 0.487. Données~:
\texttt{hf\_business.markout\_curve}.}\label{fig:hf:the-business-of-market-making:markout}
\end{omfigure}

\textbf{Pour aller plus loin.} Coter à un tick derrière les meilleurs prix et observer l'effondrement du volume~; désactiver la rémunération (exercice 7)~; donner au Quoter
un prix juste au lieu du mid (chapitre 2) et une raison d'incliner ses cotations (chapitre 3).

\section{Construction~: le banc d'essai et l'économie unitaire}\label{bld:hf:the-business-of-market-making:harness}

\begin{build}
\bfield{Rôle} Une interface unique pour chaque stratégie de market making du livre, de sorte qu'une stratégie écrite une seule fois tourne sur le marché synthétique du livre~7 aujourd'hui
et sur le simulateur de bourse du livre~10 dès qu'il sera disponible~; et l'arithmétique de la capture, des coûts et de l'échelle.
\bfield{Interface} \texttt{firm.mmharness}~: un \texttt{Quoter} avec \texttt{on\_start}, \texttt{on\_market(ctx, t, top)}, \texttt{on\_fill(ctx, fill)}
et, facultativement, \texttt{on\_timer}~; \texttt{ctx.send}, \texttt{cancel}, \texttt{take}, \texttt{quote}, \texttt{external}, \texttt{position},
\texttt{cash}, \texttt{working()}~; \texttt{run\_tape(quoter, cfg, latency, fees) -> Result} avec \texttt{fills}, \texttt{pnl(ref)},
\texttt{decompose(H, ref)} et \texttt{markouts(horizons, ref)}, \texttt{ref} étant le mid ou le prix efficient. \texttt{firm.mmecon}~:
\texttt{capture\_bp}, \texttt{profit}, \texttt{breakeven\_volume}, \texttt{operating\_leverage}, \texttt{cost\_shares}.
\bfield{Règles} Le Quoter ne voit jamais le prix efficient. Sa vue du carnet exclut ses propres ordres. Un ordre annulé peut encore être exécuté pendant que
l'annulation voyage. Même configuration et même graine, mêmes exécutions.
\bfield{Tests d'acceptation} \texttt{code/firm/mmharness/tests/}~: la décomposition se somme au P\&L marqué par rapport aux deux références~; la limite de position
tient à un ordre près en vol~; les messages restent inférieurs à un dixième de ceux du marché~; une exécution agressive paie les frais du preneur~; les exécutions sont déterministes.
\texttt{code/firm/mmecon/tests/}~: le seuil de rentabilité et le levier opérationnel par rapport à leurs définitions.
\bfield{Pour aller plus loin} Un adaptateur vers \texttt{firm.exchsim} (le livre~10, chapitre~26) avec un test de parité avec celui-ci~; des timers en temps de marché.
\end{build}

\begin{omsources}
\item Virtu Financial, Inc., Form 10-K for the year ended 31 December 2025 (SEC EDGAR).
\item Flow Traders, «~4Q and FY 2025 Results~», 12 février 2026.
\item A.~J.~Menkveld, «~High frequency trading and the new market makers~», \emph{Journal of Financial Markets} 16(4), 2013.
\item M.~Baron, J.~Brogaard, B.~Hagströmer et A.~Kirilenko, «~Risk and return in high-frequency trading~», \emph{Journal of Financial and
Quantitative Analysis} 54(3), 2019.
\end{omsources}

\section{Exercices}

\begin{exercise}[$\star$]\label{exo:hf:the-business-of-market-making:1}
Le revenu net de trading ajusté de Virtu s'élève à 2\,145.3~millions de dollars sur 248.5 jours de trading. Quel est-il par jour~?
\end{exercise}

\begin{exercise}[$\star$]\label{exo:hf:the-business-of-market-making:2}
Flow Traders a gagné 485.8~millions d'euros de revenu net de trading sur 1\,940~milliards d'euros de valeur échangée en produits négociés en bourse. Exprimer le ratio en
points de base et expliquer pourquoi il surestime la capture sur ces produits.
\end{exercise}

\begin{exercise}[$\star$]\label{exo:hf:the-business-of-market-making:3}
Une entreprise capture 0.5 point de base de la valeur échangée, supporte 0.2 point de base de coûts variables et a 40\,000~\$ de coûts fixes par jour. Quelle est
sa valeur quotidienne échangée au seuil de rentabilité~?
\end{exercise}

\begin{exercise}[$\star\star$]\label{exo:hf:the-business-of-market-making:4}
La même entreprise échange 2~milliards de dollars par jour. Calculer son profit, son levier opérationnel, et son profit après une baisse de 10~\% du volume.
\end{exercise}

\begin{exercise}[$\star\star$]\label{exo:hf:the-business-of-market-making:5}
Par rapport au mid, la capture de spread du Quoter symétrique est de $+0.487$ cent par action~; par rapport au prix efficient, elle est de $-0.173$. Expliquer comment
les deux peuvent être vraies.
\end{exercise}

\begin{exercise}[$\star\star$]\label{exo:hf:the-business-of-market-making:6}
D'après la \cref{fig:hf:the-business-of-market-making:markout}, entre quels horizons le mark-out par rapport au mid change-t-il de signe, et ce que cela
dit de la durée pendant laquelle le Quoter peut garder une position~?
\end{exercise}

\begin{exercise}[$\star\star\star$]\label{exo:hf:the-business-of-market-making:7}
\emph{Programmation.} Relancer les dix heures sans aucun frais (\texttt{Fees(0, 0)}). Quel est le P\&L total par action, et pourquoi rien d'autre ne change-t-il~?
\end{exercise}

\begin{exercise}[$\star\star\star$]\label{exo:hf:the-business-of-market-making:8}
\emph{Trouver la faille.} «~Nous avons fait tourner notre nouveau quoter pendant une heure dans le simulateur, il a rapporté 61~\$. Ça marche, augmentons la taille.~»
\end{exercise}

\section{Problème~: une fraction de centime}

\begin{problem}[{Problème du week-end --- une fraction de centime}]\label{pb:hf:the-business-of-market-making:1}
Une fondatrice veut lancer un \omterm{def:hf:the-business-of-market-making:emm}{teneur de marché électronique}. On dispose des données publiques, d'un simulateur et d'un tableur.

\medskip\noindent\textbf{Partie I --- Les données publiques.}
\begin{enumerate}
\item Définir un \omterm{def:hf:the-business-of-market-making:emm}{teneur de marché électronique} et dire ce qu'il vend.
\item Donner le revenu net de trading ajusté de Virtu en 2025 par jour, et la part qu'en représente son segment de market making.
\item Additionner les lignes de coûts fixes de Virtu et les exprimer par jour de trading.
\item Calculer le revenu net de trading ajusté de Virtu et la rémunération par employé.
\item Calculer le revenu net de trading de Flow Traders par unité de valeur échangée en produits négociés en bourse, ses coûts fixes par employé et sa marge d'EBITDA.
\end{enumerate}

\medskip\noindent\textbf{Partie II --- Économie unitaire.}
\begin{enumerate}[resume]
\item Définir le \omterm{def:hf:the-business-of-market-making:capture}{taux de capture du revenu} et écrire le profit en fonction du volume quotidien.
\item Déduire le volume au seuil de rentabilité et le levier opérationnel.
\item Pour la petite entreprise (0.5 et 0.2 points de base, 40\,000~\$ par jour), donner le volume au seuil de rentabilité, ainsi que le profit et le levier à 2~milliards de dollars par jour.
\item Quel est son profit sur 250 jours de trading à 2~milliards de dollars par jour~?
\item Les coûts fixes passent à 50\,000~\$ par jour. Quel est le nouveau volume au seuil de rentabilité~?
\item La capture tombe à 0.4 point de base. Quel est alors le volume au seuil de rentabilité, et le profit à 2~milliards de dollars~?
\end{enumerate}

\medskip\noindent\textbf{Partie III --- Le premier Quoter.}
\begin{enumerate}[resume]
\item Énoncer l'identité du P\&L de la \cref{prop:hf:the-business-of-market-making:split} et la démontrer.
\item Pourquoi un Quoter doit-il coter hors du carnet, sans ses propres ordres~?
\item Donner le volume horaire du Quoter, sa part de marché et ses messages.
\item Donner son P\&L par action par rapport au mid, composante par composante.
\item Pourquoi sa capture de spread est-elle négative par rapport au prix efficient~?
\item Quelle est la précision de la moyenne sur dix heures du total horaire~?
\end{enumerate}

\medskip\noindent\textbf{Partie IV --- Le verdict.}
\begin{enumerate}[resume]
\item Qu'ont établi Menkveld, Baron et leurs coauteurs, et qu'impliquent-ils pour ce Quoter~?
\item Énoncer le \emph{résultat nommé}~: le volume quotidien au seuil de rentabilité de la petite entreprise et son levier opérationnel à 2~milliards de dollars par jour.
\item En deux phrases, que la fondatrice doit-elle construire en premier~: davantage de capture, davantage de volume ou des coûts fixes plus bas~?
\end{enumerate}
\end{problem}

\section{Questions d'entretien}

\begin{interviewq}[$\star$ \iqroles{trader}]\label{iq:hf:the-business-of-market-making:1}
Que vend un teneur de marché, et qui le paie pour cela~?
\end{interviewq}

\begin{interviewq}[$\star$ \iqroles{chercheur}]\label{iq:hf:the-business-of-market-making:2}
Décomposez le P\&L d'un teneur de marché en capture de spread, sélection adverse et inventaire. Pourquoi la décomposition dépend-elle d'un horizon~?
\end{interviewq}

\begin{interviewq}[$\star\star$ \iqroles{trader, chercheur}]\label{iq:hf:the-business-of-market-making:3}
Les exécutions de votre quoter sont rentables au moment de l'exécution et déficitaires dix secondes plus tard. Que se passe-t-il, et que tentez-vous en premier~?
\end{interviewq}

\begin{interviewq}[$\star\star$ \iqroles{développeur}]\label{iq:hf:the-business-of-market-making:4}
Un quoter simulé envoie un million de messages par heure et ne réalise que quelques centaines d'exécutions. Quel bogue recherchez-vous~?
\end{interviewq}

\begin{interviewq}[$\star\star$ \iqroles{risque}]\label{iq:hf:the-business-of-market-making:5}
Pourquoi le profit d'un teneur de marché est-il plus volatil que son volume~?
\end{interviewq}

\begin{interviewq}[$\star\star\star$ \iqroles{chercheur}]\label{iq:hf:the-business-of-market-making:6}
Le P\&L horaire d'une stratégie a une moyenne de $-\$19$ et un écart-type de 89~\$. Combien d'heures indépendantes faudrait-il pour distinguer une vraie moyenne de
$-\$19$ de zéro à deux erreurs types près~?
\end{interviewq}
